\section*{Limitations}

\begin{itemize}
  \item \textbf{Model scale and diversity.} Experiments limited to small models (1-1.5B parameters) and two architectures (Qwen2.5, Llama-3.2). Larger models and diverse architectures remain unexplored.

  \item \textbf{Incomplete reconstruction.} While the method achieves up to 55\% KL divergence reduction, the remaining gap indicates linear projections cannot fully reconstruct monolingual distributions, suggesting non-linear language-semantic interactions at the manifold level.

  \item \textbf{Dataset coverage.} Evaluation relies on Flores Plus and TED Talks (formal text). Conversational data, technical domains, and naturally occurring code-switching patterns are not tested. Our artificial code-switching setup produces mid-sentence switches at fixed positions, which may not reflect the morphosyntactic boundaries where natural code-switching typically occurs. Evaluation on naturalistic code-switching corpora such as the LinCE benchmark~\cite{aguilar-etal-2020-lince} is an important direction for future work.

  \item \textbf{Language typology.} Limited to Indo-European and Chinese languages. Low-resource languages, right-to-left scripts, and typologically distant languages (agglutinative, tonal) unexplored. Hindi showed the weakest improvement across both models, with KL reduction (24-27\%) but minimal gains in other distributional similarity metrics, suggesting tokenization or data balance issues.

  \item \textbf{Layer-wise steering.} Steering applied only to final layers based on empirical observation. Systematic exploration of multi-layer or adaptive steering strategies not conducted.

  \item \textbf{Evaluation metrics.} While we supplement distributional metrics (KL, JS divergence) with generation-based code-switching metrics (CSI, M-Index, I-Index), we do not evaluate semantic preservation or downstream task performance. Human evaluation of output quality remains needed.

  \item \textbf{Fluency--control trade-off and coefficient fragility.} Generation-based evaluation reveals that aggressive steering coefficients, while effective at eliminating code-switching (CSI~$< 0.01$ for English-Russian and English-Hindi), can produce repetitive or degenerate output patterns. Steering coefficients are highly fragile. Small changes in magnitude shift behavior from fluent code-switched text to degenerate repetition. The CSI-optimized coefficients cluster near the maximum tested values ($-5.0$), suggesting that direct optimization of generation metrics pushes toward an aggressive regime. Combining steering with constrained decoding or adaptive coefficient scheduling could mitigate this issue.

  \item \textbf{Steering calibration.} Optimal coefficients require manual grid search per language pair and depend on the optimization target (distributional vs.\ generation-based). Principled or adaptive selection methods not developed.

  \item \textbf{Efficiency benchmarking.} No rigorous latency profiling across hardware configurations or batch sizes to validate overhead claims for production deployment.

  \item \textbf{Bilingual, English-only steering.} Our evaluation steers exclusively toward English in bilingual settings. Multilingual steering (controlling language choice among three or more languages simultaneously) and bidirectional steering (e.g., steering from English toward Spanish) are natural extensions that we leave to future work. The multilingual PCA maps (Figures~\ref{fig:qwen_lang_map},~\ref{fig:llama_lang_map}) suggest the geometric structure supports such extensions, as multiple language clusters are simultaneously separable.
\end{itemize}